\documentclass[11pt,a4paper]{article}
\usepackage{fontspec,amsmath,amssymb,amsthm,microtype}
\setmainfont{lmroman10-regular.otf}[BoldFont=lmroman10-bold.otf,ItalicFont=lmroman10-italic.otf,BoldItalicFont=lmroman10-bolditalic.otf]
\setmonofont{lmmono10-regular.otf}[Scale=MatchLowercase]
\usepackage[margin=25mm]{geometry}
\usepackage[unicode,hidelinks]{hyperref}
\setlength{\emergencystretch}{3em}
\setlength{\parskip}{0.4em}
\setlength{\parindent}{0pt}
\allowdisplaybreaks[2]
\raggedbottom
\title{\texorpdfstring{Actual phase laws and their L\ensuremath{{}^{2}} limits}{Actual phase laws and their L² limits}}
\author{LangevinLean mathematical exposition}
\date{}
\begin{document}\maketitle
\noindent Source module: \texttt{Kinetic\allowbreak{}Langevin/\allowbreak{}Lp\allowbreak{}Phase\allowbreak{}Law.\allowbreak{}lean}. The accompanying HTML pairs this mathematical exposition with the exact Lean statements and records the scope of the formalization.


The phase variable is a pair of actual L\ensuremath{{}^{2}} classes, and its law is the pushforward of their measurable representatives. Dual characteristic functions depend continuously on this pair. This identifies a limiting stationary phase law from bounded-test estimates without incorrectly treating the usual maximum product norm as a Hilbert norm.


\section{\texorpdfstring{Verify measurability of actual phase representatives}{Verify measurability of actual phase representatives}}


Each L\ensuremath{{}^{2}} class has a fixed strongly measurable representative. Taking the pair of its velocity and position representatives gives a measurable phase map. Lean uses the two strong-measurability results and the measurable product constructor, making the subsequent pushforward an actual measure.



{\raggedright\noindent Lean declarations: \texttt{measurable\_\allowbreak{}lp\allowbreak{}Phase\allowbreak{}Coe}.\par}

\section{\texorpdfstring{Reduce phase characteristic functions to one real L\ensuremath{{}^{2}} variable}{Reduce phase characteristic functions to one real L² variable}}


For a continuous real linear functional \(L\) on phase space, let \(L_v(v)=L(v,0)\) and \(L_x(x)=L(0,x)\). The real L\ensuremath{{}^{2}} variable \(R=L_v(V)+L_x(X)\) represents \(L(V(\omega),X(\omega))\) almost everywhere. Its law therefore has scalar characteristic function at frequency one equal to the dual characteristic function of the phase law. The map \((V,X)\mapsto R\) is continuous, since the two continuous linear maps lift continuously to L\ensuremath{{}^{2}} and addition is continuous. The existing scalar characteristic-function continuity theorem proves the desired phase continuity. Lean explicitly intersects the representative events for both lifted linear maps and their sum, proves the coordinate decomposition by linearity, rewrites the actual pushforward through \(L\), and composes the two continuous maps.



{\raggedright\noindent Lean declarations: \texttt{continuous\_\allowbreak{}char\allowbreak{}Fun\allowbreak{}Dual\_\allowbreak{}map\_\allowbreak{}lp\allowbreak{}Phase}.\par}

\section{\texorpdfstring{Identify the law of an actual L\ensuremath{{}^{2}} limit from vanishing weak errors}{Identify the law of an actual L² limit from vanishing weak errors}}


Suppose actual phase pairs \(S_n\) converge in the product of L\ensuremath{{}^{2}} spaces to \(S\), and their bounded complex-test errors relative to a fixed finite measure \(\pi\) are eventually bounded by \(\delta_n\to0\). For each dual functional use the norm-one test \(F(z)=e^{iLz}\). The error squeeze gives convergence of the approximating characteristic functions to that of \(\pi\); the continuity theorem gives convergence of the same sequence to the characteristic function of the actual limiting phase law. Uniqueness of limits equates them, and uniqueness of finite measures from their dual characteristic functions gives \(P\circ S^{-1}=\pi\). Lean rewrites the actual phase expectation by the measurable pushforward formula and invokes the checked scalar squeeze and measure-uniqueness theorems. The weak estimates must be proved by the concrete numerical and probability layers; stationary output is not assumed.



{\raggedright\noindent Lean declarations: \texttt{measure\_\allowbreak{}map\_\allowbreak{}lp\allowbreak{}Phase\_\allowbreak{}eq\_\allowbreak{}of\_\allowbreak{}bounded\_\allowbreak{}test\_\allowbreak{}error}.\par}

\end{document}
